% Author wording: attachment 5d87e492-ea14-4fe2-922b-d6922a14a7c0.
% Figure directions are rendered as diagrams; editorial closing note omitted.
\section{Preliminary calculations}
[Note: this part is quite long, if you are interested in the \hyperref[sec:experiments]{experiment result}, you can skip this section] \\

Before running the experiments, let’s first break down what's the time/memory complexity for each method.

to recap:
\begin{itemize}
  \item Method A performs a "query-document dot product" operation for every document inside the opened clusters.
  \item Method B avoids some "query-document dot product" operations by filtering with Bitplanes first.
  \item Method C avoids starting with a large document set by looking up a specific prefix, then broadening only when needed.
\end{itemize}

So for each method, I mainly want to understand two things:

\begin{itemize}
  \item What additional work do Methods B and C introduce?
  \item How many "query-document dot product" operations can Methods B and C potentially save?
\end{itemize}

Let's use this as variables:

\begin{center}
\begin{minipage}{\linewidth}
\small
\renewcommand{\arraystretch}{1.2}
\begin{tabular}{@{}p{.26\linewidth}p{.68\linewidth}@{}}
\toprule
Symbol & Meaning\\
\midrule
\(N\) & Total number of documents\\
\(C\) & Total number of clusters\\
\(P\) & Number of clusters opened for one query (P is for "probes")\\
\(d\) & Number of binary dimensions\\
\(K\) & Number of results we want to return\\
\(n_c\) & Number of documents inside cluster \(c\)\\
\(F_A,F_B,F_C\) & Number of documents that undergo a "query-document dot product" operation under A, B and C\\
\bottomrule
\end{tabular}
\end{minipage}
\end{center}


If the clusters are roughly balanced, one cluster contains about

\[
n \approx \frac{N}{C}
\]

documents.

So if we open \(P\) clusters, we expose roughly

\[
\frac{PN}{C}
\]

documents.




\subsection{Query-time work}


Let's try to imagine this using example. Suppose one opened cluster contains three documents:

\begin{center}
\begin{minipage}{\linewidth}
\small
\renewcommand{\arraystretch}{1.2}
\begin{tabular}{@{}p{.16\linewidth}p{.42\linewidth}p{.34\linewidth}@{}}
\toprule
ID & Binary document & Full score value\\
\midrule
0 & \texttt{10100} & 0.1\\
1 & \texttt{11001} & 2.5\\
2 & \texttt{11000} & 1.9\\
\bottomrule
\end{tabular}
\end{minipage}
\end{center}


The query is

\[
q=[0.7,\ 0.5,\ -0.4,\ -0.6,\ 0.3].
\]

From the previous section, we know that the ideal binary document is

\texttt{11001}

and suppose we want the best

\[
K=2
\]

documents.

For this example:

\[
N=3,\qquad C=P=1,\qquad d=5.
\]

*note: the actual experiments later contain hundreds of thousands or millions of documents.



\subsubsection{Method A: scan the opened clusters}

Method A is the easiest one to understand.

Once cluster selection (prefilter) has decided which clusters to open, \textbf{for every document inside those clusters, we will do a "query-document dot product" operation}.

\Needspace{4\baselineskip}
\textbf{Step 1: how many "query-document dot product" operations do we need?}\par

In our example, the only opened cluster contains:

\texttt{[ID 0, ID 1, ID 2]}

so:

\[
F_A=3.
\]

More generally, let \(O\) be the set of clusters opened for one query. Then:

\[
\boxed{
F_A=\sum_{c\in O}n_c
}
\]

is the number of documents for which A performs "query-document dot product" operations.

If the clusters are roughly balanced:

\[
F_A\approx\frac{PN}{C}.
\]

So for A, this is already the most important number: \textbf{how many documents did cluster selection leave for us to search?}


\begin{figure}[htbp]
\centering
\begin{tikzpicture}

\node[scan,text width=16mm] (q) at (0,0) {Query};
\node[scan,text width=24mm] (p) at (29mm,0) {Select $P$\\clusters};
\node[scan,text width=29mm] (rows) at (65mm,0) {$F_A$ documents\\\texttt{10100}\\\texttt{11001}\\\texttt{11000}};
\node[scan,text width=29mm] (score) at (104mm,0) {"query-document dot product" operation\\for every document};
\node[scan,text width=16mm] (top) at (138mm,0) {Top $K$};
\draw[arr] (q)--(p); \draw[arr] (p)--(rows); \draw[arr] (rows)--(score); \draw[arr] (score)--(top);

\end{tikzpicture}
\caption{Method A performs no additional filtering after cluster selection. Every document inside the opened clusters undergoes a "query-document dot product" operation.}
\end{figure}





\Needspace{4\baselineskip}
\textbf{Step 2: what does one "query-document dot product" operation cost?}\par

The scorer used here follows the four-sign lookup-table approach described in Exa’s article.

Instead of multiplying every query value with every document bit individually, the query is divided into groups (or subvectors) of four.

For our five-dimensional example:

\begin{center}\begingroup\small\ttfamily\begin{tabular}{l}
[0.7, 0.5, -0.4, -0.6]\\ \relax [0.3, 0, 0, 0]
\end{tabular}\endgroup\end{center}

A four-bit group has

\[
2^4=16
\]

possible binary patterns.

So this example needs

\[
2\times16=32
\]

precomputed table entries.

In general, the number of groups is

\[
G=\left\lceil\frac d4\right\rceil.
\]

One "query-document dot product" operation therefore reads \(G\) table entries.

For our example:

\[
G=2.
\]

Three documents therefore require:

\[
3\times2=6
\]

table lookups.

In general, the scoring work grows with

\[
F_AG.
\]

Since \(G\) grows with \(d\),

\[
T_{\text{score},A}=O(F_Ad).
\]

For a fixed binary dimension, this is basically linear in the number of documents for which we perform "query-document dot product" operations.



\Needspace{4\baselineskip}
\textbf{Step 3: keep only the best K}\par

As score values are produced, we only need to keep the best \(K\).

For our example:

\begin{center}\begin{tikzpicture}
\node[scan,text width=30mm] (a) at (0,0) {ID 0: 0.1};
\node[scan,text width=30mm] (b) at (46mm,0) {ID 1: 2.5\\ID 0: 0.1};
\node[scan,text width=30mm] (c) at (92mm,0) {ID 1: 2.5\\ID 2: 1.9};
\draw[arr](a)--(b);\draw[arr](b)--(c);
\end{tikzpicture}\end{center}

Using a bounded heap, maintaining the best \(K\) results takes at most approximately

\[
O(F_A\log(K+1)).
\]



\Needspace{4\baselineskip}
\textbf{Total cost for A}\par

Let \(T_{\text{base}}\) contain the work before "query-document dot product" operations, such as cluster selection and lookup-table preparation.

Let \(a_A\) be the average measured time needed to fully process one document, including its "query-document dot product" operation and the top-\(K\) update.

Then:

\[
\boxed{
T_A\approx T_{\text{base}}+F_Aa_A.
}
\]

For roughly balanced clusters:

\[
\boxed{
T_A\approx
T_{\text{base}}
+
\frac{PN}{C}a_A.
}
\]

The main tradeoff is straightforward.

Using more clusters usually makes each cluster smaller, so each opened cluster contains fewer documents. However, we may then need to open more clusters to maintain the same recall.

\Needspace{6\baselineskip}
\begin{center}\setlength{\fboxsep}{9pt}\fcolorbox{Ink!25}{Soft}{\parbox{.91\linewidth}{\textbf{TL;DR, Method A:} A’s cost is mainly determined by how many documents cluster selection leaves behind. Fewer opened documents means fewer "query-document dot product" operations, but opening too few clusters can hurt recall.}}\end{center}



\subsubsection{Method B: Bitplanes}

Bitplanes try to reduce the number of documents that undergo a "query-document dot product" operation.

Instead of immediately performing a "query-document dot product" operation for every document inside the opened clusters, B first uses Bitplanes to filter the documents.

So B has two main costs:

\[
\boxed{
\text{Bitplane work}
+
\text{"query-document dot product" operations for the documents that survive}
}
\]

The whole idea only makes sense if the "query-document dot product" operations we avoid are more expensive than the Bitplane work required to avoid them.



\Needspace{4\baselineskip}
\textbf{Step 1: inspect the important query bits first}\par

B uses the same cluster selection and the same lookup table as A.

The additional step is to order the query bits by decreasing \(|q_i|\).

For

\[
q=[0.7,\ 0.5,\ -0.4,\ -0.6,\ 0.3],
\]

the order is:

\begin{center}\begin{tikzpicture}
\foreach \i/\bit/\weight in {0/1/0.7,1/4/0.6,2/2/0.5,3/3/0.4,4/5/0.3} {
\node[branch,text width=19mm] (b\i) at (\i*29mm,0) {Bit \bit\\(\weight)};
}
\foreach \i/\j in {0/1,1/2,2/3,3/4} {\draw[arr](b\i)--(b\j);}
\end{tikzpicture}\end{center}

We inspect the larger values first because mismatching them hurts the score value more.

Sorting these positions costs

\[
O(d\log d).
\]

We also calculate the ideal score value:

\[
Q=\sum_i|q_i|=2.5.
\]


\begin{figure}[htbp]
\centering
\begin{tikzpicture}

\node[scan,text width=30mm] (p) at (0,0) {Select $P$ clusters};
\node[branch,text width=56mm] (o) at (54mm,0) {Order bits by $|q_i|$\\$1\to4\to2\to3\to5$};
\node[branch,text width=32mm] (b) at (111mm,0) {Bitplane search};
\draw[arr] (p)--(o); \draw[arr] (o)--(b);

\end{tikzpicture}
\caption{Bitplanes keep the same cluster-selection stage as A, but add query-dependent filtering before the full scorer.}
\end{figure}





\Needspace{4\baselineskip}
\textbf{Step 2: represent the documents using a mask}\par

Inside one cluster, B keeps a bitmap showing which document positions are still active.

If cluster \(c\) contains \(n_c\) documents, its mask needs

\[
\boxed{
W_c=
\left\lceil
\frac{n_c}{64}
\right\rceil
}
\]

64-bit machine words.

Our three-document example fits inside one word.

The starting mask is:

\texttt{111}

which means that all three documents are still active.



\Needspace{4\baselineskip}
\textbf{Step 3: split using the query bits}\par

The first query position is bit 1, and the query prefers \texttt{1}.

All three documents match:

\begin{center}\begingroup\small\ttfamily\begin{tabular}{l}
111 AND 111 = 111
\end{tabular}\endgroup\end{center}

so nothing is removed.

The next position is bit 4. The query prefers \texttt{0}, and again all three documents match.

So the mask is still:

\texttt{111}.

The next position is bit 2.

Its Bitplane is:

\texttt{011}.

Therefore:

\begin{center}\begingroup\small\ttfamily\begin{tabular}{l}
111 AND 011 = 011
\end{tabular}\endgroup\end{center}

The preferred branch now contains IDs 1 and 2.

The alternative branch is:

\texttt{100}

which contains ID 0.

Because ID 0 disagrees with the query at bit 2, this branch has mismatch penalty:

\[
p=|q_2|=0.5.
\]


\begin{figure}[htbp]
\centering
\begin{tikzpicture}

\node[branch,text width=43mm] (root) at (0,0) {\texttt{111}: IDs 0,1,2\\mask width = 1 word};
\node[branch,text width=43mm] (first) at (0,-22mm) {\texttt{111}\\mask width = 1 word};
\node[branch,text width=43mm] (fourth) at (0,-44mm) {\texttt{111}\\mask width = 1 word};
\node[branch,text width=47mm] (yes) at (-32mm,-69mm) {\texttt{011}: IDs 1,2; $p=0$\\mask width = 1 word};
\node[box,draw=gray!65,fill=gray!7,dashed,text width=47mm] (no) at (32mm,-69mm) {\texttt{100}: ID 0; $p=0.5$\\mask width = 1 word};
\draw[arr] (root)--node[right,font=\small]{Bit 1}(first);
\draw[arr] (first)--node[right,font=\small]{Bit 4}(fourth);
\draw[arr] (fourth)--node[above left,font=\small]{Bit 2 = 1}(yes);
\draw[arr,dashed] (fourth)--node[above right,font=\small]{Bit 2 = 0}(no);

\end{tikzpicture}
\caption{The first two bit tests remove no documents. The third separates ID 0 into another branch, but that branch is not discarded yet.}
\end{figure}





\Needspace{4\baselineskip}
\textbf{Step 4: why Bitplane filtering is not O(1), it's not free! :(}\par

At first, a bitwise AND seems like it's O(1).

But actually it's not. Example -> If one cluster contains:

\[
6,400
\]

documents.

Then, for every one mask we need:

\[
6400/64=100
\]

machine words.

So one split does not perform one AND, as we need to loop across roughly 100 word positions.

There is another important detail: \textbf{the mask does not physically get shorter when fewer documents remain active.}

For example, suppose only 80 documents survive after several splits. But in this case, the mask can still be 100 words wide.


\begin{figure}[htbp]
\centering
\begin{tikzpicture}

\node[anchor=west,font=\small\bfseries] at (0,10mm) {Initial mask: 6,400 active documents};
\node[anchor=east,font=\small] at (140mm,10mm) {100 words};
\node[anchor=west,font=\small\bfseries] at (0,-18mm) {Later mask: 80 active documents};
\node[anchor=east,font=\small] at (140mm,-18mm) {100 words};
\foreach \i in {0,...,24} {
  \draw[draw=Branch!55,fill=Branch!70] (\i*5.6mm,0) rectangle ++(5.1mm,5mm);
  \draw[draw=Branch!35,fill=Branch!5] (\i*5.6mm,-28mm) rectangle ++(5.1mm,5mm);
}
\fill[Branch] (0,-28mm) rectangle (1.75mm,-23mm);
\draw[<->,draw=Ink!65] (0,-35mm)--node[below,font=\small]{Same physical width}(140mm,-35mm);

\end{tikzpicture}
\caption{The number of active documents can become much smaller while the physical bitmap stays the same size.}
\end{figure}

\pagebreak  





\Needspace{4\baselineskip}
\textbf{Step 5: count the Bitplane work}\par

Let \(J_c\) be the number of times B splits a mask in cluster \(c\).

Each split scans \(W_c\) words.

So the total amount of split-mask work is:

\[
\boxed{
V_{\text{split}}
=
\sum_{c\in O}J_cW_c.
}
\]

For roughly balanced clusters, if B performs an average of \(S\) splits per opened cluster:

\[
V_{\text{split}}
\approx
PS
\left\lceil
\frac{N}{64C}
\right\rceil.
\]



\Needspace{4\baselineskip}
\textbf{Step 6: eventually stop splitting and perform "query-document dot product" operations}\par

After a branch contains few enough documents, we treat this as a \textbf{leaf}, and in this case, method B will perform "query-document dot product" operations for the documents inside it.

Suppose our leaf size is two.

The preferred branch:

\texttt{011}

contains exactly IDs 1 and 2.

So B performs these "query-document dot product" operations:

\begin{center}\begingroup\small\ttfamily\begin{tabular}{l}
ID 1 \(\rightarrow\) 2.5\\ \relax ID 2 \(\rightarrow\) 1.9
\end{tabular}\endgroup\end{center}

Therefore:

\[
F_B=2.
\]

*in this case, if we do a comparison, method A would have performed three "query-document dot product" operations.



\Needspace{4\baselineskip}
\textbf{another little cost: we still need to read leaf masks}\par

Before B can extract the surviving document positions, it still has to scan the leaf mask. Let \(H_c\) be the number of leaf masks read in cluster \(c\).

\[
\boxed{
V_{\text{leaf}}
=
\sum_{c\in O}H_cW_c.
}
\]




\Needspace{4\baselineskip}
\textbf{Step 7: the un-explored branch that is still in the queue}\par

Remember when we consider 'branching' to catch all cases because we are worried that we might miss the ideal case due to branching? Imagine that: ID 0 is still waiting to be explored.

We can calculate the known mismatch penalty, which is:

\[
p=0.5.
\]

However, we can take a look that the best possible score value inside that branch is:

\[
Q-2p
=
2.5-2(0.5)
=
1.5.
\]

Meanwhile, our current top two score values are already:

\[
2.5,\qquad1.9.
\]

Since:

\[
1.5<1.9,
\]

We can make conclusion: nothing inside the waiting branch can enter the current top two. So the good news is that, we can safely skip it!


\begin{figure}[htbp]
\centering
\begin{tikzpicture}

\node[box,draw=gray!65,fill=gray!7,dashed,text width=49mm,minimum height=24mm] (wait) at (0,0) {Waiting branch\\penalty = 0.5\\best possible score value = 1.5};
\node[branch,text width=35mm,minimum height=24mm] (best) at (82mm,0) {Current top-2\\2.5\\1.9};
\draw[arr] (wait)--node[above,font=\small]{$1.5<1.9$}(best);
\node[font=\small\bfseries,text=Branch] at (41mm,-23mm) {$\rightarrow$ skip branch};

\end{tikzpicture}
\caption{Once the best possible score value of a branch falls below the current top-\(K\) threshold, the whole branch can be skipped.}
\end{figure}





\Needspace{4\baselineskip}
\textbf{Total cost for B}\par

B therefore pays for:

\begin{itemize}\setlength{\itemsep}{1pt}
\item ordering the query bits;
\item reading split masks;
\item reading leaf masks;
\item maintaining the waiting branches;
\item performing \(F_B\) "query-document dot product" operations.
\end{itemize}

Let:

\begin{center}
\begin{minipage}{\linewidth}
\small
\renewcommand{\arraystretch}{1.2}
\begin{tabular}{@{}p{.26\linewidth}p{.68\linewidth}@{}}
\toprule
Symbol & Meaning\\
\midrule
\(b_s\) & Average time for one split-mask word\\
\(b_l\) & Average time for one leaf-mask word\\
\(T_{\text{order}}\) & Ordering query bits and preparing query weights\\
\(T_{\text{queue}}\) & Managing branches, masks and the priority queue\\
\(a_B\) & Average time to fully process one surviving document\\
\bottomrule
\end{tabular}
\end{minipage}
\end{center}


Then:

\[
\boxed{\begin{aligned}T_B\approx{}&T_{\text{base}}+T_{\text{order}}+b_sV_{\text{split}}+b_lV_{\text{leaf}}\\&+F_Ba_B+T_{\text{queue}}.\end{aligned}}
\]

The key condition is:

\[
\boxed{
\text{"query-document dot product" work avoided}
>
\text{Bitplane work added}.
}
\]

A smaller leaf size may reduce \(F_B\), because B keeps filtering longer. However, that also means more splits and more bitmap work.

A larger node budget can also improve recall, but again it allows B to search more branches.

\Needspace{6\baselineskip}
\begin{center}\setlength{\fboxsep}{9pt}\fcolorbox{Ink!25}{Soft}{\parbox{.91\linewidth}{\textbf{TL;DR, Bitplanes:} It is fast when a few Bitplane splits can filter out most documents before the "query-document dot products" operations. If the search needs many splits, or the masks remain wide, the bitmap work grows and can cancel out the savings from scoring fewer documents.}}\end{center}



\subsubsection{Method C: Backward Walk}
Now let's look at the second idea.

Backward Walk approaches the problem from almost the opposite direction.

Method A starts with every document inside the opened clusters.

Method B also starts with all of them, but progressively removes documents.

Backward Walk instead asks:

\textbf{What if we start from the most specific binary prefix first, and only broaden the search when we need more documents?}

So C mainly pays for two things:
\begin{itemize}\setlength{\itemsep}{1pt}
\item looking up different prefixes;
\item performing "query-document dot product" operations for the documents exposed by those prefixes.
\end{itemize}

\Needspace{4\baselineskip}
\textbf{Step 1: store the documents in prefix order}\par
Using the same example, let's sort the three documents by their binary key:
\begin{center}
\begin{tabular}{@{}rlr@{}}\toprule
Position & Binary key & ID\\\midrule
0 & \texttt{10100} & 0\\
1 & \texttt{11000} & 2\\
2 & \texttt{11001} & 1\\\bottomrule
\end{tabular}
\end{center}
Meanwhile, the ideal query pattern is \texttt{11001}.

Because the rows are sorted, documents with the same prefix will appear next to each other.

That means we can use binary search to quickly find the range corresponding to a particular prefix.

\Needspace{4\baselineskip}
\textbf{Step 2: start from the most specific prefix}\par
Let's first use all five bits: \texttt{11001}.

Two binary searches find the beginning and the end of this range.

In our example, the result is $[2,3)$, which contains only ID 1.

So we perform one "query-document dot product" operation: ID 1 $\to$ 2.5.

However, remember that we want the top two results. Right now we only have one document.

So unfortunately, we cannot stop yet. We need to broaden the prefix.

\Needspace{4\baselineskip}
\textbf{Step 3: remove one constraint}\par
We now go from \texttt{11001} to \texttt{1100*}.

The \texttt{*} means that we no longer care whether the last bit is \texttt{0} or \texttt{1}.

The new range becomes $[1,3)$, which contains:
\begin{itemize}\setlength{\itemsep}{1pt}
\item ID 2: \texttt{11000};
\item ID 1: \texttt{11001}.
\end{itemize}
Now, ID 1 was already scored at the previous depth.

So the good news is that we do \textbf{not} need to score it again.

The only new document is ID 2: ID 2 $\to$ 1.9.

Now we have two results.

\begin{figure}[htbp]\centering
\begin{tikzpicture}
\node[box,draw=gray!60,fill=gray!10,text width=35mm] (outside) at (-44mm,0) {ID 0\\\texttt{10100}};
\node[keys,fill=Keys!12,text width=35mm] (added) at (0,0) {ID 2\\\texttt{11000}};
\node[keys,fill=Keys!38,text width=35mm] (old) at (44mm,0) {ID 1\\\texttt{11001}};
\node[font=\small,text=Keys] (label) at (0,15mm) {new document};
\draw[arr](label)--(added);
\draw[Keys,thick] (25mm,-9mm)--(25mm,-12mm)--(63mm,-12mm)--(63mm,-9mm);
\node[font=\small,text=Keys] at (44mm,-16mm) {\texttt{11001}};
\draw[Keys,thick] (-19mm,-21mm)--(-19mm,-24mm)--(63mm,-24mm)--(63mm,-21mm);
\node[font=\small,text=Keys] at (22mm,-28mm) {\texttt{1100*}};
\end{tikzpicture}
\caption{The broader prefix contains the previous range. ID 1 is already scored, so only the newly included document, ID 2, needs a "query-document dot product" operation.}
\end{figure}

\Needspace{4\baselineskip}
\textbf{Step 4: how much prefix lookup work do we need?}\par
Suppose we start at prefix depth $x$, and eventually stop at depth $\ell$.

Then the number of depths we visit is:
\[
L=x-\ell+1.
\]
At every positive depth, we need two binary searches for every opened cluster:
\begin{itemize}\setlength{\itemsep}{1pt}
\item one to find where the prefix starts;
\item one to find where it ends.
\end{itemize}
So if we open $P$ clusters and visit $L_+$ positive prefix depths, we perform $2PL_+$ binary searches.

Depth zero uses the known whole-cluster range and needs no binary search.

Each binary search over cluster $c$ costs roughly:
\[
O(\log(n_c+1)).
\]
So the total lookup work is approximately:
\[
2L_+\sum_{c\in O}\log_2(n_c+1).
\]
Again, the formula looks more complicated than the idea.

The idea is simply:

\textbf{Backward Walk cost grows with how many prefix depths we visit and how many clusters we search at each depth.}

\Needspace{4\baselineskip}
\textbf{Step 5: how many documents do we eventually score?}\par
Let $F_C$ be the total number of distinct documents that become visible before we stop.

Because each broader prefix contains the previous prefix, a document that was already scored does not need another "query-document dot product" operation.

So C's document-scoring work is simply:
\[
F_Ca_C.
\]
Here $a_C$ is the average time for a "query-document dot product" operation and its top-$K$ update.

This means the main question for Backward Walk is actually:

\textbf{How broad does the prefix need to become before we have enough useful documents?}

If we can stop while the prefix is still very specific, $F_C$ can be very small.

If we need to keep broadening all the way to the root, $F_C$ eventually becomes all documents in the opened clusters.

\Needspace{4\baselineskip}
\textbf{Step 6: how quickly does the prefix grow?}\par
Let's use a simple model just to build intuition.

Assume, temporarily, that every bit is independent and has a 50\% chance of being \texttt{0} or \texttt{1}.

Then the probability that a random document matches an $\ell$-bit prefix is:
\[
2^{-\ell}.
\]
So if we have $n$ documents, the expected number matching that prefix is:
\[
\frac{n}{2^\ell}.
\]
For a collection of one million documents:
\begin{center}
\begin{tabular}{@{}rr@{}}\toprule
Prefix depth & Expected documents\\\midrule
16 & about 15\\
15 & about 31\\
14 & about 61\\
13 & about 122\\
12 & about 244\\\bottomrule
\end{tabular}
\end{center}
So something interesting happens here.

We only remove \textbf{one bit} each time, but the number of matching documents approximately \textbf{doubles}.

\begin{figure}[htbp]\centering
\begin{tikzpicture}
\node[anchor=east,font=\small]at(-3mm,0){16 bits $\to$ about 15};
\node[anchor=east,font=\small]at(-3mm,-9mm){15 bits $\to$ about 31};
\node[anchor=east,font=\small]at(-3mm,-18mm){14 bits $\to$ about 61};
\node[anchor=east,font=\small]at(-3mm,-27mm){13 bits $\to$ about 122};
\node[anchor=east,font=\small]at(-3mm,-36mm){12 bits $\to$ about 244};
\draw[Keys,fill=Keys!25](0,-2.5mm)rectangle(6mm,2.5mm);
\draw[Keys,fill=Keys!25](0,-11.5mm)rectangle(12mm,-6.5mm);
\draw[Keys,fill=Keys!25](0,-20.5mm)rectangle(24mm,-15.5mm);
\draw[Keys,fill=Keys!25](0,-29.5mm)rectangle(48mm,-24.5mm);
\draw[Keys,fill=Keys!25](0,-38.5mm)rectangle(96mm,-33.5mm);
\end{tikzpicture}
\caption{Under the independent 50/50-bit model, removing one prefix bit doubles the expected matching count. The labels are rounded; they describe document counts, not recall.}
\end{figure}

However, we should remember that this is only an intuition.

Real embedding bits can be correlated, and once the documents have already been divided into clusters, the distribution can be even further from a perfect 50/50 split.

Most importantly:

\textbf{This only estimates how many documents we might find. It does not tell us whether those documents are actually the correct top results.}

\Needspace{4\baselineskip}
\textbf{Step 7: when should Backward Walk stop?}\par
There are two ways to stop.

\textbf{Approximate stopping}\par
The simplest option is to stop after we have found enough documents.

For example:
\begin{center}
\begin{tabular}{@{}rl@{}}
16 bits & $\to$ 12 documents\\
15 bits & $\to$ 31 documents\\
14 bits & $\to$ 66 documents\\
13 bits & $\to$ 136 documents
\end{tabular}
\end{center}
If our target is 100 documents, we could stop at 13 bits.

This is simple and cheap. However, there is one problem.

Finding 100 documents does not prove that the best unseen document is worse than the ones we found.

So this stopping rule is approximate, and we need to measure its recall experimentally.

\textbf{Exact stopping}\par
Fortunately, we can also use the scoring equation to derive a safe bound.

Suppose we have already searched all documents matching the first $\ell$ query signs inside the opened clusters.

Any unseen document in those clusters must disagree with the query on at least one of those first $\ell$ positions.

The cheapest possible mismatch is:
\[
m_\ell=\min_{1\le i\le\ell}|q_i|.
\]
Therefore, the best possible score of an unseen document is:
\[
U_\ell=Q-2m_\ell.
\]
Meanwhile, suppose we have already kept $K$ documents and the worst document currently inside our top-$K$ has score $\tau$.

If:
\[
\tau>U_\ell,
\]
then even the \textbf{best possible unseen document} cannot enter our top-$K$.

So the good news is that we can stop safely. The implementation also keeps its tie-breaking and floating-point rounding checks.

If this never happens, we can keep broadening the prefix until depth zero.

At depth zero, we have exposed every document inside the opened clusters, so Backward Walk becomes equivalent to Method A locally.

\Needspace{4\baselineskip}
\textbf{Total cost for Backward Walk}\par
Backward Walk therefore pays for:
\begin{itemize}\setlength{\itemsep}{1pt}
\item building the query prefix;
\item performing prefix lookups;
\item updating the current prefix ranges;
\item performing "query-document dot product" operations for $F_C$ exposed documents.
\end{itemize}
We can summarize this as:
\[
\begin{aligned}
T_C\approx{}&T_{\text{base}}+T_{\text{key}}
+2L_+t_b\sum_{c\in O}\log_2(n_c+1)\\
&+F_Ca_C+T_{\text{ranges}}.
\end{aligned}
\]
$T_{\text{key}}$ covers the query key and optional bound preparation; $t_b$ is time per binary-search comparison; $T_{\text{ranges}}$ covers range updates and stopping checks.

Again, the exact symbols are not the most important part.

The important condition is:
\[
\boxed{\text{"query-document dot product" work avoided}>\text{prefix lookup work added}}
\]
If C can stop at a fairly deep prefix, then $F_C\ll F_A$, and we may save a lot of scoring work.

However, if C has to broaden the prefix all the way to depth zero, $F_C=F_A$, then C ends up doing all of A's "query-document dot product" operations \textbf{plus} the additional prefix lookups.

\Needspace{6\baselineskip}
\begin{center}\setlength{\fboxsep}{9pt}\fcolorbox{Ink!25}{Soft}{\parbox{.91\linewidth}{\textbf{TL;DR, Backward Walk:} It is fast when a specific prefix already contains the documents we need. Otherwise, it must keep relaxing the prefix and expanding the searched range; if this goes too far, Backward Walk approaches a full scan while also paying the extra binary-search cost for each new prefix range.}}\end{center}

\subsection{Memory complexity}
So far, we have only talked about query time.

However, B and C also need additional index structures, so we should also check how much memory they require.

There are two kinds of memory that matter here:
\begin{itemize}\setlength{\itemsep}{1pt}
\item \textbf{persistent index memory}, which remains loaded between queries;
\item \textbf{temporary query memory}, which only exists while one query is running.
\end{itemize}

\subsubsection{Method A memory}
Method A stores:
\begin{itemize}\setlength{\itemsep}{1pt}
\item the packed binary document;
\item the document ID;
\item the cluster/routing information.
\end{itemize}
Using 64-bit words, one binary document needs
\[
8\left\lceil\frac d{64}\right\rceil
\]
bytes. The ID adds another 8 bytes.

So the document storage is:
\[
M_{A,\text{rows}}=N\left(8\left\lceil\frac d{64}\right\rceil+8\right).
\]
During a query, A also keeps the query values, the lookup table and the current top-$K$.

\subsubsection{Method B memory}
B also uses what A has been used, and additional two more things.
First, it stores the Bitplanes.

For each cluster, every binary dimension needs one bitmap spanning all document positions in that cluster.

Therefore:
\[
M_{\text{planes}}=8d\sum_{c=1}^{C}W_c.
\]
So in a rough sense, B stores the document bits twice:
\begin{itemize}\setlength{\itemsep}{1pt}
\item once arranged by document;
\item once arranged as Bitplanes.
\end{itemize}
The second extra cost is temporary query memory.

Remember that branching means we can have several masks waiting in the queue at the same time.

Each one still has the full physical width of its cluster.

If $A_c(t)$ masks belonging to cluster $c$ are alive at time $t$, then the raw mask memory is:
\[
Q_{\text{masks}}=8\max_t\sum_c A_c(t)W_c.
\]

\begin{figure}[htbp]\centering
\begin{tikzpicture}[node distance=8mm]
\node[box,text width=37mm,minimum height=31mm] (stored) {\textbf{Stored once}\\Packed documents\\IDs\\Bitplanes};
\node[box,text width=37mm,minimum height=31mm,right=of stored] (common) {\textbf{Common query}\\Query\\Lookup table\\Top $K$};
\node[branch,text width=37mm,minimum height=31mm,right=of common] (extra) {\textbf{B-specific query}\\\phantom{Mask}\\\phantom{Mask}\\\phantom{Mask}\\\phantom{...}};
\node[branch,text width=27mm,minimum height=4mm,font=\scriptsize]at([yshift=3mm]extra.center){Mask};
\node[branch,text width=27mm,minimum height=4mm,font=\scriptsize]at([yshift=-2mm]extra.center){Mask};
\node[branch,text width=27mm,minimum height=4mm,font=\scriptsize]at([yshift=-7mm]extra.center){Mask};
\node[font=\small,text=Branch]at([yshift=-12mm]extra.center){\ldots};
\end{tikzpicture}
\caption{B stores Bitplanes alongside its packed documents. During a query, several full-width masks may also remain alive in the branch queue.}
\end{figure}

% \Needspace{6\baselineskip}
% \begin{center}\setlength{\fboxsep}{9pt}\fcolorbox{Ink!25}{Soft}{\parbox{.91\linewidth}{\textbf{TL;DR, B memory:} Bitplanes need extra persistent memory because we store another arrangement of the document bits. During search, B can also temporarily hold several full-width branch masks at the same time.}}\end{center}

\subsubsection{Method C memory}
C also keeps the normal packed documents and IDs.

The current implementation additionally stores one 32-bit prefix key for every document. So:
\[
M_{\text{prefix}}=4N.
\]
The important part is that we do \textbf{not} create a new copy of the documents for every prefix depth.

For example, \texttt{11001}, \texttt{1100*}, \texttt{110**} all refer to different ranges of the \textbf{same sorted document array}.

\begin{figure}[htbp]\centering
\begin{tikzpicture}
\node[box,text width=28mm]at(-52.5mm,0){\texttt{10100}};
\node[keys,text width=28mm]at(-17.5mm,0){\texttt{11000}};
\node[keys,fill=Keys!30,text width=28mm]at(17.5mm,0){\texttt{11001}};
\node[keys,text width=28mm]at(52.5mm,0){\texttt{11010}};
\draw[Keys,thick](2.5mm,8mm)--(2.5mm,11mm)--(32.5mm,11mm)--(32.5mm,8mm);
\node[font=\small,text=Keys]at(17.5mm,15mm){\texttt{11001}};
\draw[Keys,thick](-32.5mm,20mm)--(-32.5mm,23mm)--(32.5mm,23mm)--(32.5mm,20mm);
\node[font=\small,text=Keys]at(0,27mm){\texttt{1100*}};
\draw[Keys,thick](-32.5mm,32mm)--(-32.5mm,35mm)--(67.5mm,35mm)--(67.5mm,32mm);
\node[font=\small,text=Keys]at(17.5mm,39mm){\texttt{110**}};
\node[font=\small]at(0,-11mm){One sorted array, not three copies};
\end{tikzpicture}
\caption{In this illustrative array, three prefix ranges refer to the same stored rows. Broadening the prefix does not create a new document list.}
\end{figure}

During the query, C mainly needs to remember the previous range for every opened cluster.

So compared with B, method C's temporary search state is relatively small. 

% \Needspace{6\baselineskip}
% \begin{center}\setlength{\fboxsep}{9pt}\fcolorbox{Ink!25}{Soft}{\parbox{.91\linewidth}{\textbf{TL;DR, C memory:} Backward Walk adds a small prefix key per document, but it does not need full-width masks for many waiting branches.}}\end{center}

\subsection{Recall}
Now we know roughly how much work each method performs.

However, none of those savings matter if we stop too early and return the wrong documents.

So we also need to understand where recall can be lost.

For query $q$:
\begin{itemize}\setlength{\itemsep}{1pt}
\item $G_q$ is the exact global top-$K$;
\item $O_{a,q}$ is the result returned by method $a$.
\end{itemize}
Average recall is:
\[
R_a=\frac{\sum_q|G_q\cap O_{a,q}|}{n_qK}.
\]
Here $n_q$ is the number of queries. The reference and returned results use the same score and document-ID tie rule.

For example:
\begin{center}
\begin{tabular}{@{}ll@{}}
Reference: & \texttt{[1,2]}\\
Returned: & \texttt{[1,0]}
\end{tabular}
\end{center}
means we recovered one of two: 50\%.

\Needspace{4\baselineskip}
\textbf{Where can the correct document disappear?}\par
There are actually two places.

First, cluster selection can fail to open the cluster containing the correct document.

Second, even if the correct document is inside an opened cluster, B or C may stop their local search before reaching it.

\begin{figure}[htbp]\centering
\begin{tikzpicture}[node distance=8mm]
\node[box,text width=120mm] (global) {Global exact top-$K$};
\node[box,text width=120mm,below=9mm of global] (opened) {Correct documents that reached the opened clusters};
\node[scan,text width=37mm,below left=10mm and -32mm of opened] (a) {\textbf{A}\\Scan everything\\No additional local loss};
\node[branch,text width=37mm,below=10mm of opened] (b) {\textbf{B}\\Bitplane search\\May stop early};
\node[keys,text width=37mm,below right=10mm and -32mm of opened] (c) {\textbf{C}\\Prefix search\\May stop early};
\draw[arr](global)--node[right,font=\small]{cluster selection}(opened);
\draw[arr](opened)--(a);\draw[arr](opened)--(b);\draw[arr](opened)--(c);
\end{tikzpicture}
\caption{Cluster selection can lose a correct document before local search begins. B and C may also lose a document by stopping their local search early.}
\end{figure}

\subsubsection{Recall for Method A}
Let $g_q$ be the number of global top-$K$ documents whose clusters were opened.

Because A scores every document inside those clusters, all $g_q$ of those documents are recovered locally. So:
\[
R_{A,q}=\frac{g_q}{K}.
\]
In other words:

A can lose a correct document because cluster selection missed it, but once the document is inside an opened cluster, A will not lose it during local search.

\subsubsection{Recall for Bitplanes}
Suppose $g_q$ correct documents reached the opened clusters, but B only recovers some fraction $r_{B,q}$ of them.

Then:
\[
R_{B,q}=\frac{g_q}{K}r_{B,q}.
\]
If $g_q=0$, recall is zero.

For example, imagine:
\begin{itemize}\setlength{\itemsep}{1pt}
\item 8 of the global top 10 documents are inside the opened clusters;
\item B recovers 6 of those 8.
\end{itemize}
Then:
\[
\mathrm{Recall}=\frac{8}{10}\times\frac{6}{8}=60\%.
\]
If B performs a complete branch-and-bound search using valid score bounds, then it eventually recovers all of the correct documents inside the opened clusters.

However, if we impose a finite node budget, B may stop early.

\Needspace{6\baselineskip}
\begin{center}\setlength{\fboxsep}{9pt}\fcolorbox{Ink!25}{Soft}{\parbox{.91\linewidth}{\textbf{TL;DR, B recall:} B can lose recall in two places: cluster selection can miss the document before B starts, and a limited Bitplane search can miss a document even after its cluster has been opened.}}\end{center}

\subsubsection{Recall for Backward Walk}
For Backward Walk, let $D_{q,\ell}$ be all documents exposed when the search stops at prefix depth $\ell$.

Then:
\[
R_{C,q}=\frac{|G_q\cap D_{q,\ell}|}{K}.
\]
The important thing here is that \textbf{candidate count and recall are not the same thing}.

Finding 100 documents does not mean that we found the correct 100 documents.

Consider:
\[
q=[0.1,\ 0.9].
\]
The ideal binary pattern is \texttt{11} and suppose the only documents are \texttt{10} and \texttt{01}.

Their scores are:
\[
S(10)=-0.8,\qquad S(01)=0.8.
\]
If we search prefix \texttt{1*}, we find \texttt{10}.

However, the better document is \texttt{01}, which is outside that prefix.

So even though we found one candidate, we did not find the best result.

That is why stopping after ``enough documents'' is only an approximate rule.

C is exact locally only when:
\begin{enumerate}\setlength{\itemsep}{1pt}
\item we broaden all the way to depth zero and expose every document inside the opened clusters; or
\item the score bound proves that none of the unseen documents can beat the current top-$K$.
\end{enumerate}

\Needspace{6\baselineskip}
\begin{center}\setlength{\fboxsep}{9pt}\fcolorbox{Ink!25}{Soft}{\parbox{.91\linewidth}{\textbf{TL;DR, C recall:} A large prefix range does not automatically mean high recall. Backward Walk can safely stop only when it has either searched every document or mathematically proved that the unseen documents cannot win.}}\end{center}

\subsection{Summary}
The whole calculation section can be summarized with one simple idea:
\begin{center}\small
\begin{tabularx}{\linewidth}{@{}lY@{}}\toprule
Method & What it does\\\midrule
A & Score every document inside the opened clusters\\
B & Spend Bitplane work to avoid some "query-document dot product"\\
C & Spend prefix-lookup work to avoid some "query-document dot product"\\\bottomrule
\end{tabularx}
\end{center}
So:

\textbf{A just pays for "query-document dot product". B and C deliberately add extra search work, hoping that this lets them avoid enough "query-document dot product" to make the total query faster.}

The experiment in the next section will tell us whether that actually happens on real embeddings.
